\chapter{Evaluación, resultados y discusión}
\label{cap:evaluacion}

Los capítulos anteriores han presentado la arquitectura Open CVN, su
implementación por capas y el resultado funcional que ofrece a quien la
utiliza. El presente capítulo cierra ese recorrido con una pregunta distinta:
qué evidencia respalda, de forma reproducible, cada una de las afirmaciones
técnicas realizadas hasta ahora. La evaluación no se limita a comprobar que
el sistema se ejecuta sin errores, sino que verifica, capa a capa, que cada
contrato descrito en los capítulos anteriores se cumple sobre el paquete
oficial CVN y sobre los artefactos que de él se derivan.

\section{Diseño de la evaluación}
\label{sec:evaluacion_diseno}

La evaluación se apoya en una batería de pruebas automatizadas organizada
según las mismas capas de la arquitectura presentada en el
Capítulo~\ref{cap:metodologia}, en lugar de tratarse como un bloque
monolítico. Cada nivel de evaluación comprueba una responsabilidad concreta:

\begin{itemize}
  \item \textbf{Generación estructural:} produce artefactos íntegros e
  importables.

  \item \textbf{Normalización y resolución de referencias auxiliares:}
  mantienen la trazabilidad esperada.

  \item \textbf{Política semántica:} aplica sus reglas de forma consistente
  con la evidencia disponible.

  \item \textbf{JSON Schema generado:} valida correctamente documentos
  válidos e inválidos.

  \item \textbf{Parsers e importadores:} cumplen su contrato público.

  \item \textbf{Almacenamiento y versionado:} preservan la integridad de los
  datos.

  \item \textbf{Exportación:} produce artefactos deterministas.

  \item \textbf{Flujos completos de la herramienta:} funcionan de extremo a
  extremo.
\end{itemize}

Esta batería se ejecuta localmente durante el desarrollo y, además, de forma
automática en cada contribución propuesta sobre las ramas principales del
repositorio mediante integración continua, de modo que ninguna
funcionalidad nueva se incorpora sin que el conjunto completo de garantías ya
verificadas siga cumpliéndose.

\section{Niveles de evaluación}
\label{sec:evaluacion_niveles}

La Tabla~\ref{tab:evaluacion_categorias} resume los ocho niveles de
evaluación considerados, el objetivo de cada uno y el número de pruebas
automatizadas que lo cubren actualmente.

\begin{table}[htbp]
  \centering
  \small
  \renewcommand{\arraystretch}{1.15}
  \begin{tabular}{p{0.27\textwidth} p{0.57\textwidth} r}
    \hline
    \textbf{Nivel de evaluación} & \textbf{Objetivo} & \textbf{Pruebas} \\
    \hline
    Artefactos generados & Bindings estructurales, modelos de dominio, modelo conceptual y diagramas íntegros e importables & 146 \\
    Normalización y trazabilidad & Unificación manual/árbol por código CVN y resolución de referencias auxiliares correctas & 90 \\
    Política semántica & Formas semánticas y elegibilidad de enumeración consistentes con la evidencia disponible & 76 \\
    JSON Schema y validación runtime & Validación estructural y en tiempo de ejecución de documentos válidos e inválidos & 25 \\
    Parsers e importadores & Contrato público de análisis, importación JSON/XML/PDF y casos de error & 63 \\
    Almacenamiento y versiones & Persistencia SQLite, currículo maestro y versiones derivadas & 25 \\
    Exportación & Generación determinista de LaTeX y PDF & 19 \\
    Flujos de extremo a extremo & Pipeline completo y flujo completo de la herramienta local & 44 \\
    \hline
    \textbf{Total} & & \textbf{488} \\
    \hline
  \end{tabular}
  \caption{Categorías de pruebas automatizadas y objetivo de cada una.}
  \label{tab:evaluacion_categorias}
\end{table}

Esta organización por niveles permite además delimitar responsabilidades:
una regresión detectada en el nivel de política semántica, por ejemplo, se
diagnostica sin necesidad de examinar la herramienta local o los mecanismos
de exportación, porque cada nivel verifica un contrato independiente de los
demás.

\section{Comando de verificación y resultados}
\label{sec:evaluacion_comando}

La verificación completa del sistema se ejecuta con un único comando:

\begin{code}[htbp]
\begin{lstlisting}[language=, basicstyle=\small\ttfamily, breaklines=true]
uv run pytest -n auto tests
\end{lstlisting}
\caption{Comando principal de verificación del sistema.}
\label{cod:evaluacion_comando}
\end{code}

Este comando ejecuta en paralelo las 488 pruebas de la Tabla~\ref{tab:evaluacion_categorias}
sobre el paquete oficial CVN vigente en el repositorio, sin datos simulados
ni resultados precalculados. La ejecución más reciente sobre esta versión
del documento superó las 488 pruebas sin ningún fallo, en 692,80 segundos
(11 minutos y 32 segundos) de ejecución paralela. El mismo comando es el que
ejecuta la integración continua ante cualquier contribución propuesta sobre
las ramas principales del repositorio, por lo que el resultado no depende
del entorno de quien lo ejecuta.

\section{Discusión de garantías y limitaciones}
\label{sec:evaluacion_discusion}

Superar la batería de pruebas certifica que el sistema se comporta según lo
descrito en los capítulos anteriores, pero no todas las partes del sistema
ofrecen el mismo tipo de garantía. Conviene distinguir, en particular, las
limitaciones que proceden del propio paquete oficial CVN de las que
proceden de una decisión deliberada de alcance de este Trabajo Fin de Grado.

Las limitaciones que proceden del paquete oficial no son atribuibles al
sistema desarrollado, sino a inconsistencias o a la propia naturaleza del
material fuente:

\begin{itemize}
  \item \textbf{Discrepancia XML/XSD del modelo en árbol.} El XML canónico y
  su esquema de validación no coinciden por completo, tal y como se
  documentó en el Capítulo~\ref{cap:ecosistema_cvn}.

  \item \textbf{Referencia \texttt{CVN\_AGENCY\_C} sin resolver.} El paquete
  no ofrece una tabla equivalente localizable para esta referencia del
  manual.

  \item \textbf{Deriva de empaquetado entre familias auxiliares.} Obliga a
  un mapeo consciente del paquete en lugar de un convenio de nombres fijo.
\end{itemize}

El sistema no oculta estas circunstancias: las registra explícitamente y las
trata mediante evidencia en lugar de forzar una interpretación no
justificada.

Las limitaciones que proceden del alcance del propio trabajo son, en cambio,
decisiones deliberadas que delimitan hasta dónde llega la propuesta actual:

\begin{itemize}
  \item \textbf{Importación de CVN XML.} Es un mapeo semántico parcial y
  creciente, no un conversor completo de cualquier documento CVN XML real.

  \item \textbf{JSON Schema generado.} No puede expresar toda regla
  semántica del dominio ni todas las relaciones conceptuales curadas.

  \item \textbf{Importación asistida por LLM.} Es opcional, y su resultado
  no se considera autoritativo hasta que se revisa.

  \item \textbf{Generación de PDF.} Depende de la disponibilidad de un motor
  LaTeX en el entorno de ejecución.
\end{itemize}

La Tabla~\ref{tab:evaluacion_garantias} resume estas garantías junto con la
limitación conocida que las acompaña.

\begin{table}[htbp]
  \centering
  \small
  \renewcommand{\arraystretch}{1.15}
  \begin{tabular}{p{0.4\textwidth} p{0.44\textwidth}}
    \hline
    \textbf{Garantía verificada} & \textbf{Limitación conocida asociada} \\
    \hline
    Generación estructural reproducible desde los esquemas XSD oficiales & Algunas construcciones XSD (\textit{choice}, cardinalidades mínimas) no se preservan de forma estricta en los bindings generados \\
    Normalización y resolución de referencias auxiliares con evidencia trazable & Una referencia del paquete analizado permanece sin resolver (\texttt{CVN\_AGENCY\_C}) \\
    Política semántica evaluada a partir de evidencia concreta, no de listas fijas & Algunas tablas compactas quedan como revisión pendiente en lugar de enumeración cerrada o catálogo abierto definitivo \\
    Validación estructural y en tiempo de ejecución determinista del formato Open CVN JSON & El JSON Schema no puede expresar toda relación conceptual curada ni toda regla semántica del dominio \\
    Importación de CVN XML con trazabilidad y diagnóstico explícito & Mapeo semántico parcial, creciente sobre los casos ya cubiertos; no es un conversor completo \\
    Importación asistida por LLM opcional y validada antes de almacenarse & Depende de un proveedor externo; el resultado no es autoritativo sin revisión \\
    Exportación determinista a LaTeX y PDF & La generación de PDF depende de un motor LaTeX disponible en el entorno \\
    \hline
  \end{tabular}
  \caption{Garantías del sistema frente a limitaciones conocidas.}
  \label{tab:evaluacion_garantias}
\end{table}

En conjunto, el sistema desarrollado ofrece garantías fuertes y
reproducibles allí donde el proceso es determinista: la generación
estructural, la normalización, la resolución de referencias con evidencia
concreta, la política semántica y la validación del formato Open CVN JSON.
Ofrece garantías parciales y explícitamente documentadas allí donde el
propio dominio o el alcance del trabajo lo exigen: la importación de CVN
XML, la importación asistida por LLM y la generación de PDF. Ninguno de
estos resultados debe interpretarse como una conversión completa del
ecosistema CVN ni como una validación semántica total del dominio
curricular; ambas interpretaciones excederían lo que la evidencia reunida en
este capítulo permite afirmar. El Capítulo~\ref{cap:conclusiones} retoma esta distinción para
valorar el grado de cumplimiento de los objetivos planteados en el
Capítulo~\ref{cap:introduccion} y para señalar las líneas de trabajo futuro
que se derivan de las limitaciones aquí documentadas.
